\subsection{Automi come grafi}
A partire da un automa finito $(T,q_0,F)$ se ne può dare una descrizione grafica nel modo seguente: si consideri il grafo i cui nodi sono gli stati $Q$, e gli archi etichettati con elementi di $A$ che collegano due stati $(q_i,q_j)$ se esiste la transizione $T(q_i,a)=q_j$.
inoltre avremo una funzione che ci indica lo stato iniziale
\[
\texttt{init}(G_{(T,q_0,F)}) = q_0
\]
(di solito indicata da un arco entrante nel nodo iniziale $q_0$) ed una che distingue gli stati accettanti (risp. non accettanti)
\[
\texttt{acc}(G_{(T,q_0,F)}) = F \qquad (\text{risp. } \texttt{ref}(G_{(T,q_0,F)}) = Q \setminus F)
\]
(di solito si indicano gli stati accettanti con un doppio cerchio, mentre tutti gli altri nodi sono considerati non accettanti).
Per essere più formali:
\begin{defi}[Grafo associato ad un automa]
Dato un automa $(T,q_0,F)$ definito su un alafabeto $A$ e insieme di stati $Q$ definiamo il suo grafo associato $G=(V,E,L)$ come segue:
\begin{itemize}
    \item l'insieme dei vertici è
    \[
    V:=Q;
    \]
    \item l'insieme dei lati è
    \[
    E:=\{(q_i,a,q_j)\in Q\times A\times Q \mid T(q_i,a)=q_j\};
    \]
    \item la funzione di etichettatura è
    \[
    L:E\to A, \qquad L(q_i,a,q_j)=a.
    \]
\end{itemize}
\end{defi}

Ad esempio il grafo associato all'automa dell'Esempio~\ref{ex: parole lunghezze pari}, è rappresentato in Figura~\ref{figura esempio automa per parole pari}.
\begin{figure}[ht]
\centering
\begin{tikzpicture}[
    ->, >=stealth', shorten >=1pt, auto, node distance=3.5cm,
    every state/.style={circle, draw, minimum size=1cm}
]
    \node[state, accepting] (q0) {$q_0$};
    \node[state] (q1) [right of=q0] {$q_1$};

    \draw[->] (-1.2,0.6) -- (q0);
    \path (q0) edge [bend left=30] node {$a$} (q1)
          (q1) edge [bend left=30] node {$a$} (q0);
\end{tikzpicture}
\caption{Automa dell'Esempio \ref{ex: parole lunghezze pari}, notare che $a$ è un qualsiasi elemento nell'alfabeto $A=\{0,1\}$.}
\label{figura esempio automa per parole pari}
\end{figure}
\begin{figure}[h]
\centering
\begin{tikzpicture}[
    ->, >=stealth', shorten >=1pt, auto, node distance=4cm,
    every state/.style={circle, draw, minimum size=1cm}
]
    \node[state, accepting] (q0) {$q_0$};
    \node[state] (q1) [right of=q0] {$q_1$};

    \draw[->] (-1.2,0.6) -- (q0);

    \path (q0) edge [loop above] node {$0$} (q0)
          (q1) edge [loop above] node {$0$} (q1)
          (q0) edge [bend left=25] node {$1$} (q1)
          (q1) edge [bend left=25] node {$1$} (q0);
\end{tikzpicture}
\caption{In modo simile all'automa dell'Esempio~\ref{ex: parole lunghezze pari} si può dare l'automa che decide l'insieme di tutte le parole con un numero pari di 1.}
\label{fig:parita}
\end{figure}

% =========================================================
% FIGURA 3
% =========================================================
\begin{figure}[ht]
\centering
\begin{tikzpicture}[
    ->, >=stealth', shorten >=1pt, auto, node distance=3.5cm,
    every state/.style={circle, draw, minimum size=1cm}
]
    \node[state] (q0) {$q_0$};
    \node[state, accepting] (q1) [right of=q0] {$q_1$};
    \node[state] (q2) [below right=2cm and 0.3cm of q0] {$q_2$};

    \draw[->] (-1.2,0.6) -- (q0);

    \path (q0) edge [bend left=15] node {$0$} (q1)
          (q1) edge [loop above] node {$1$} (q1)
          (q0) edge [bend right=15] node {$1$} (q2)
          (q1) edge [bend left=15] node {$0$} (q2)
          (q2) edge [loop below, min distance=10mm, in=200, out=250] node[below left] {$0$} (q2)
          (q2) edge [loop below, min distance=10mm, in=290, out=340] node[below right] {$1$} (q2);
\end{tikzpicture}
\caption{Automa che decide l'insieme $X=\{01^n \mid n\geq 0\}$.}
\label{fig:01n}
\end{figure}

\begin{comment}
\subsection{Rappresentazione come matrice di adiacenza}

Possiamo usare $M_E$ per indicare la matrice di adiacenza del grafo $G=(V,E)$ definita come 
\[
(M_E)_{ij} =
\begin{cases}
1 & \text{se } (v_i,v_j) \in E \\
0 & \text{altrimenti.}
\end{cases}
\]
Ad esempio, in Figura~\ref{fig:grafo7}, abbiamo un grafo orientato e la corrispondente matrice di adiacenza.
È sensato, che ad ogni insieme $E$ corrisponda un'unica matrice una volta che è fissato un ordinamento degli elementi di $V$. Utilizzando le  operazioni aritmetiche  date dal semianello dei booleani $(\mathbb{B}, \vee, \wedge)$:

\[
\mathbb{B} = \{0,1\} \qquad
\begin{array}{c|cc}
a \vee b & a & b \\
\hline
0 & 0 & 0 \\
1 & 1 & 0 \\
1 & 0 & 1 \\
1 & 1 & 1
\end{array}
\qquad
\begin{array}{c|cc}
a \wedge b & a & b \\
\hline
0 & 0 & 0 \\
0 & 1 & 0 \\
0 & 0 & 1 \\
1 & 1 & 1
\end{array}
\]
possiamo definire il prodotto seguente
\[
A=(a_{ij})\in\mathbb{B}^{m\times n},
\qquad
B=(b_{ij})\in\mathbb{B}^{n\times p},
\]
\[
(AB)_{ij}
:=
\bigvee_{k=1}^{n}
\left(a_{ik}\wedge b_{kj}\right),
\qquad
1\le i\le m,\;1\le j\le p.
\]
Usando il prodotto definito sopra,
la matrice $M_E^2$ ha un 1 in posizione $i,j$ se esiste un cammino di lunghezza esattamente 2 che va dal vertice $v_i$ al vertice $v_j$, questo equivale a dire che esiste (almeno) un vertice $w$ tale che $(v_i,w),(w,v_j) \in E$.
Inoltre, si può provare per induzione su $k$ il seguente fatto:
\begin{prop}
       $(M_E^k)_{ij}=1$ se e solo se esiste un cammino di lunghezza esattamente $k$ da $v_i$ a $v_j$.
\end{prop}

\begin{rema}\textcolor{red}{controlla}
 La chiusura riflessiva e transitiva di $G$ si può ottenere come limite delle somme parziali delle potenze successive di $M_E^k$:
\[
M_E^* = \sum_{i\geq 0} M_E^i
\]

Si noti come in questo caso la somma infinita può essere limitata da un intero $L$ che certamente esiste e che corrisponde alla massima lunghezza del cammino più breve che connette una coppia di nodi (per ogni coppia). Una volta calcolata la potenza corrispondente a tale lunghezza, il valore della somma non sarà modificata ulteriormente:
\[
M_E^* = \sum_{i=0}^{L} M_E^i.
\]
\end{rema}

\end{comment}

\subsection{Matrici di adiacenza per grafi di automi finiti}
Consideriamo un automa 
$(T,q_0,F)$ definito su un alafabeto 
$A$ e insieme di stati $Q$.
Possiamo sempre considerare il monoide $M=(A^*,A^\emptyset, \cdot)$ 
e l'anello (che e/' anche un semianello) $\mathbb N$.
A questo punto possiamo definire l'insieme delle serie formali
\[
\mathbb N [M]
:=
\left\{
    \sum_{w \in A^*}n_w w \, | \, n_w \in \mathbb N 
\right\}
\]
Ora possiamo considerare la matrice di adiacenza $Z$ dell'automa,
che non sarà nient'altro che una matrice cha nelle entrate elementi di $\mathbb N[M]$.
Come sappiamo, fare la potenza $k$-esima di una matrice di adiacenza con etichette,
fornisce nell'entrata $(i,j)$
tutti e i soli cammini di lunghezza $k$ che collegano il vertice
$i$ al vertice $j$ (nel caso degli automi, dallo stato $q_i$ allo stato $q_j$).
Vediamo un esempio:
\begin{exam}
    Consideriamo l'automa che decide l'inisieme delle parole di lunghezza pari. 
    Questo e/' rappresentato dal seguente grafo etichettato.
\begin{center}
\begin{tikzpicture}[
    ->, >=stealth', shorten >=1pt, auto, node distance=4cm,
    every state/.style={circle, draw, minimum size=1cm}
]
    % Stati (q0 ha il doppio cerchio con 'accepting')
    \node[state, accepting] (even) {$q_0$};
    \node[state] (odd) [right of=even] {$q_1$};

    % Freccia iniziale esterna
    \draw[->] (-1.2,0.6) -- (even);

    % Transizioni SOPRA (da q0 a q1)
    \path (even) edge [bend left=45] node[above] {$1$} (odd);
    \path (even) edge [bend left=15] node[above] {$0$} (odd);

    % Transizioni SOTTO (da q1 a q0)
    \path (odd) edge [bend left=45] node[below] {$1$} (even);
    \path (odd) edge [bend left=15] node[below] {$0$} (even);
\end{tikzpicture}
\end{center}
La sua matrice di adiacenza e/'
\[
Z=
\begin{pmatrix}
   \bot & 0 + 1
    \\
    0+1 & \bot
\end{pmatrix}
\]
Infatti, gli archi $0$ e $1$ collegano i vertici 
$q_0$ e $q_1$ in entrambe le direzioni.
Inoltre, visto che non ci sono loop, 
sulla diagonale si inserisce convenzionalmente la serie formale vuota, 
indicata con $\bot$.
A questo punto vogliamo trovare i cammini di lunghezza 2 del grafo.
Per farlo basta calcolare la potenza seconda di $Z$:
\begin{align*}
Z^2&=
\begin{pmatrix}
    \bot \bot + (0+1)(0+1) & \bot(0+1) + (0+1)\bot
    \\[6pt]
    (0+1)\bot + \bot(0+1) & (0+1)(0+1) + \bot \bot
\end{pmatrix}
\\ &=
\begin{pmatrix}
    00+ 01+ 10+11 & \bot + \bot
    \\[6pt]
    \bot + \bot &  00+ 01+ 10+11
\end{pmatrix}
\\ &=
\begin{pmatrix}
    00+ 01+ 10+11 &  \bot
    \\[6pt]
    \bot &  00+ 01+ 10+11
\end{pmatrix}
\end{align*}
Per verificare le operazioni effettuate con $\bot$, 
basta usare la definizione presente in (\ref{def:serie-formali}).
Dalla matrice possiamo quindi dedurre che
esistono 4 cammini che collegano $q_0$ e $q_1$ con loro stessi,
e nessuno che collega $q_0$ con $q_1$ o viceversa.
\end{exam}

A questo punto ci vogliamo porre la seguente domanda: 
si può dedurre dala matrice di adiacenza tutti i cammini che uniscono lo stato iniziale
ad uno stato accettante dell'automa?
In altre parole:
dalla matrice di adiacenza si può dedurre il calcolo che l'automa compie per decidere un insieme?
La risposta e' si.
In partiolare, e/' anche possibile dedurre l'intero insieme $X$ da decidere.
Prima di affrontare il caso generale
analizziamo questo problema nel caso dell'esempio precedente.
\begin{exam}
    In questo caso esiste un solo stato accettante,
    che coincide esattamente con lo stato iniziale.
    Quindi,
    in una qualsiasi potenza di $Z$,
    dovremmo considerare l'entrata $(1,1)$
    per trovare le parole accettate dall'automa.
    Più precisamente, per ogni $k$ naturale, nell'entrata $(1,1)$
    della matrice $Z^k$
    troviamo le uniche e sole parole
    di lunghezza $k$ accettate dall'automa.
\end{exam}
Lavoriamo ora al caso generale. Se con $Z$ indichiamo la matrice di adiacenza di un automa,
possiamo definire
definiamo
\[
E(Z)
:=
\sum_{k=0}^\infty Z^k
\]
dove $Z^0$ e/' la matrice cha ha sulla diagonale la parola vuota $A^\emptyset$
e nelle altre entrate la serie vuota $\bot$.
Assumendo che lo stato iniziale sia $q_0$
e che corrisponda alla prima riga della matrice $Z$,
allora sulla prima riga di $E(Z)$
troveremo le parole che l'automa accetta nelle colonne corrispondenti
agli stati accettanti.
\begin{exam}
    Facendo ancora un paio di potenze di $Z$, 
    si può dedurre la seguente formula:
    \[
    Z^k=
    \left\{
    \begin{matrix}
        \begin{pmatrix}
            (0+1)^k & \bot
            \\
            \bot & (0+1)^k
        \end{pmatrix}
        ,& \text{se } k \text{ e/' pari}
        \\
        \\
        \begin{pmatrix}
            \bot & (0+1)^k
            \\
            (0+1)^k & \bot
        \end{pmatrix}
        ,& \text{se } k \text{ e/' dispari}
    \end{matrix}
    \right.
    \]
    quindi abbiamo
    \begin{align*}
    E(Z)&=
    \begin{pmatrix}
        A^\emptyset & \bot
        \\
        \bot & A^\emptyset
    \end{pmatrix}
    +
    \begin{pmatrix}
            \bot & (0+1)
            \\
            (0+1) & \bot
    \end{pmatrix}
    +
    \begin{pmatrix}
            (0+1)^2 & \bot
            \\
            \bot & (0+1)^2
        \end{pmatrix}
    + \cdots
    \\&=
    \begin{pmatrix}
        A^\emptyset + \sum_{k =1}^{\infty}(0+1)^{2k} & \bot + \sum_{k =0}^{\infty}(0+1)^{2k+1}
        \\
        \bot + \sum_{k =0}^{\infty}(0+1)^{2k+1} &  A^\emptyset + \sum_{k =1}^{\infty}(0+1)^{2k} 
    \end{pmatrix}
    \\&=
    \begin{pmatrix}
        A^\emptyset + \sum_{k =1}^{\infty}(0+1)^{2k} & \sum_{k =0}^{\infty}(0+1)^{2k+1}
        \\
        \sum_{k =0}^{\infty}(0+1)^{2k+1} &  A^\emptyset + \sum_{k =1}^{\infty}(0+1)^{2k} 
    \end{pmatrix}
    \end{align*}
    Per quanto osservato prima,
    nel posto $(1,1)$ della matrice
    $E(Z)$ troviamo le uniche e solo parole accettate dall'automa.
    \textcolor{red}{Non coverrebbe scrivre E(Z) a partire con k=1? perhce ha priori non e/' vero che la parola vuota e/' decidibile}
\end{exam}
Usando un po' di algebra lineare,
e/' possibile esplicitare esattamente l'inisieme delle parole accettate.
Infatti, fi possiamo definire i seguenti vettori:
\[
V_0(i)=
\left\{
\begin{matrix}
    A^\emptyset,  & q_i=q_0 
    \\
    \bot, & q_i\neq q_0 
\end{matrix}
\right.
\]
\[
V_F(i)=
\left\{
\begin{matrix}
    A^\emptyset,  & q_i \in F
    \\
    \bot, & q_i\notin F
\end{matrix}
\right.
\]
E calcolando il prodotto righe per colonne
\[
V_0 \cdot E(Z) \cdot V_F^t
\]
otteniamo la serie formale che ha come addendi esattamente le parole accettate da $X$.
La formula appena enunciata prende il nome di \textit{formula dell'esecuzione}.
\footnote{Ricordiamo che $\bot$ e/' l'elemento neutro additivo, 
che $A^\emptyset$ e/' l'elemento neutro moltiplicativo,
e che $\bot \cdot w = \bot$ per ogni serie formale $w \in \mathbb N[M]$.
}
\begin{exam}
    Nel caso dell'esempio precedente abbiamo
    \[
    V_0=(A^\emptyset, \bot)
    \]
    \[
    V_F=(A^\emptyset, \bot)
    \]
    Facendo il prodotto abbiamo
    \begin{align*}
        V_0 \cdot E(Z) \cdot V_F^t
        &=
        \begin{pmatrix}
              A^\emptyset & \bot    
        \end{pmatrix}
        \cdot 
        \begin{pmatrix}
        A^\emptyset + \sum_{k =1}^{\infty}(0+1)^{2k} & \sum_{k =0}^{\infty}(0+1)^{2k+1}
        \\
        \sum_{k =0}^{\infty}(0+1)^{2k+1} &  A^\emptyset + \sum_{k =1}^{\infty}(0+1)^{2k} 
    \end{pmatrix}
        \cdot
        \begin{pmatrix}
            A^\emptyset
            \\
            \bot
        \end{pmatrix}
        \\
        &=
        \begin{pmatrix}
            A^\emptyset + \sum_{k =1}^{\infty}(0+1)^{2k}
            &
            \sum_{k =0}^{\infty}(0+1)^{2k+1}
        \end{pmatrix}
        \cdot
        \begin{pmatrix}
            A^\emptyset
            \\
            \bot
        \end{pmatrix}
        \\&=
         A^\emptyset + \sum_{k =1}^{\infty}(0+1)^{2k}
    \end{align*}
    Che e/' esattamente quanto abbiamo osservato precedentemente.
\end{exam}